Sources : études citées, présentées dans le complément « Les élasticités estimées sur données françaises » ; chaque chiffre, avec sa page et son tableau, figure dans donnees/elasticites\_france.csv.\par Lecture : chaque ligne mesure la même chose, de combien un revenu déclaré augmente quand la part que le contribuable garde de l'euro marginal de ce revenu augmente ; 0,2 signifie +2 \% de revenu pour +10 \% de part gardée. Les estimations diffèrent encore par la population et l'horizon : variations d'une année sur l'autre, sauf pour Bach et al., qui comparent plusieurs années avant et après chaque réforme parmi les seuls assujettis à l'ISF. Les élasticités d'autres grandeurs (patrimoine déclaré, résidence, fraude, dons, participation) ne se placent pas sur cette échelle. La bande grise marque la fourchette de 0,12 à 0,40 que Saez, Slemrod et Giertz (2012) retiennent pour le revenu imposable aux États-Unis ; les flèches, des valeurs hors de l'échelle.
